\section{Evidence comparisons and causal boundaries}\label{sec:C12:evidence}
ContextPilot separates the editing toolset from the learning rule. Its tools include structured memory, notes, planning, and reversible context offload. Training uses snapshots at editing boundaries, masks previously optimized tokens, and allocates continuation sampling around consequential editing actions~\cite{pan2026contextpilot}. Its action-level feedback addresses the fact that final task correctness is an imperfect label for every earlier edit. This does not make the feedback causal ground truth: continuation outcomes still depend on the downstream policy and environment. A comparison must separate the expanded toolset, supervised initialization, branch allocation, and advantage estimator rather than attributing their joint effect to one component.

CompactionRL jointly trains execution and summary generation with token-level normalization. Its cross-segment credit correction is approximate, not exact full-trajectory GAE~\cite{li2026compactionrl}. For segment $s$ with $n_s\geq1$ optimized tokens, discount $\gamma\in(0,1]$ and trace parameter $\lambda\in(0,1]$, the method scales its locally computed advantages by the number of later optimized tokens:
\begin{equation}
 \widehat A_{s,i}=(\gamma\lambda)^{N_{>s}}A^{\mathrm{loc}}_{s,i},
 \qquad N_{>s}=\sum_{j>s}n_j.
 \label{eq:C12:compactionrl-position}
\end{equation}
This restores the terminal-reward discount distance when the shared reward is placed at each independently optimized segment end. It does not recursively propagate all later TD residuals into earlier segments. The authors trade that reconstruction for parallel segment processing; a reproduction must distinguish reward-distance correction from full-trajectory GAE. Their coding evaluation uses a specified 200-task SWE-bench Verified subset, rather than the full benchmark. Successful joint optimization does not establish preservation of arbitrary future constraints.

TRACE v2 evaluates candidate summaries through paired closed-loop continuations from the same AppWorld state while keeping the agent and proposer weights frozen~\cite{cumulative_a_trace_v2_20261029}. Its local verifier counts the union of blocked actions and canonicalized repeated tool calls over a fixed horizon; this does not detect every silent state corruption. Preferences between summaries reuse the same PRE continuations. The proposer receives the downstream system prompt, compressor input, and better--worse summary pairs, but not rollout traces or the verifier decomposition. Five candidate compression templates are then selected on development-set success in both repeated runs, not success in either run; the selected template is frozen for test evaluation. This separates boundary-local supervision, prompt selection, and held-out testing. The evidence remains concentrated in AppWorld, and restoration fidelity and verifier scope limit the causal interpretation.

PAIR likewise uses counterfactual continuations to identify damaging individual compressions and adapts sections of a structured prompt~\cite{min2026pair}. It distinguishes repeated-run reliability from finding one successful attempt. Its offline continuation search and restored-state access are part of the adaptation cost. These requirements can be reasonable in an experimental environment while unavailable for an external business tool with irreversible side effects. TRACE and PAIR share research lineage, so their agreement is evidence of a developing mechanism rather than two independent replications.

These mechanisms expose a concrete disagreement with purely semantic compression objectives. A summary can retain entities and facts yet obscure whether an action already happened, whether a plan is obsolete, or whether the task is complete. In that case the failure is execution-state mislocalization, not necessarily absence of a fact. Evaluating summary similarity would miss the directional role of action-observation history. The appropriate response is to preserve validated progress and effect receipts, then test whether the agent stops, resumes, and avoids repeating completed work. This is a synthesis of the mechanisms above, not a new measured result.

ACM provides author-reported evidence that explicit context actions can improve its evaluated tasks while reducing peak active context~\cite{corpus260723809}. The same study also illustrates why resource metrics must remain separate: in its reported BrowseComp-Plus comparison, tool calls increase from 19.5 to 46.2 per episode. This observation does not negate a context benefit; it limits the inference that smaller context necessarily means lower total latency or expense. These results have not been independently reproduced for this monograph.

Comparing learned compaction with deterministic summarization requires matching more than model identity. A learned controller may have seen environment-specific trajectories, received a larger training budget, or benefited from a retriever tuned to the evaluation interface. An ablation that removes the controller but also removes access to raw history changes two mechanisms simultaneously. Preserve the same storage and retrieval interfaces while varying the decision policy if the claim concerns control rather than memory availability.

Author-reported improvements across several benchmarks can support broader relevance, but shared infrastructure creates dependence between results. The same backbone, tool wrapper, answer extraction, and evaluator may appear across tasks. A cross-domain average does not turn those results into independent replications. Report task-specific outcomes and explain which components are shared. If one task improves and another regresses, the average should not hide the mechanism-specific disagreement.

The key comparison axes are behavior preservation, resource accounting, reversibility, training consistency, and implementation transparency. Full-history methods preserve observations but can suffer distraction and expense. Summaries reduce active input but can lose qualifications. Retrieval retains scalable access but can miss implicit relevance. Editable contexts increase flexibility but complicate cache and rollout semantics. These are structural tradeoffs, not a ranking without workload assumptions.

Evidence should also identify what remains unavailable. If a paper describes prompts but releases no execution adapter, code-level reproduction is incomplete. If code exists but requires a private checkpoint or undisclosed environment, the visible repository does not close that gap. A reproducibility claim should identify model weights, data, tool fixtures, edit policy, and evaluation version. Successful compilation of a harness is not successful reproduction of its scientific results.

The practical implication is to choose the simplest context-control mechanism that satisfies measured task and authority constraints. A sophisticated learned editor can be appropriate when decision value changes rapidly and enough training feedback exists. A deterministic structured state can be preferable when obligations are explicit and failures are costly. The evidence needed to choose between them is a matched workload study, not the elegance of the editing interface.

An October 2026 alternative is transparent-compact for the Pi coding agent. The inspected paper evaluates version 0.8.0, retaining user messages and assistant prose verbatim while masking tool outputs and reasoning, explicitly marking removed material, and providing delegated recovery~\cite{woolford2026transparent}. Its recovery action is \path{context_lookup}, with bounded fallback access. The guarantees are budget-limited: when retained prose reaches the configured cap, truncation begins. It therefore does not establish indefinite verbatim memory. Its next-turn probes concern commitments and planted facts with file and shell tools removed to isolate conversational memory. They do not establish complete software-engineering task success.

This method exposes a different comparison axis from learned abstraction. It preserves what was said while declining to treat that prose as verified truth. A retained assistant claim can still be wrong, and a removed tool result can require recovery before the claim is reused. The working view gains transparency about missing evidence at the cost of carrying more conversational text. Whether that tradeoff improves a task depends on the importance of exact commitments, the recoverability of observations, and the model's tendency to act confidently on an ambiguous summary.

The pinned CLM implementation supplies a concrete interface boundary. In \path{clm/clm_harness/context_env/env.py}, \path{ContextEnv._apply_edit} round-trips candidate messages through \path{parse_back}, counts tokens, and invokes \path{edit_gate.may_grow} before replacement~\cite{repo_clm2026}. The gate constrains growth; a shrinking edit that remains over budget can still be accepted with a warning, so this helper alone is not an absolute budget-admission guard. In \path{clm/clm_harness/context_utils/context_string.py}, \path{parse_back} restores the protected prefix but maps every editable non-assistant role to user text and does not reconstruct structured tool calls. Legal serialization therefore does not preserve all original role and effect provenance. Evaluate those changes separately from semantic commitment preservation and from the stricter proposed admission contract in Eq.~\eqref{eq:C12:context-admissibility}.
